\secciongrande{Mecanismos de ejecución de código en un navegador web la arquitectura del navegador}

\resumen{
  La interacción en un navegador inicia con la solicitud de un recurso y finaliza con su visualización, adaptándose a prioridades como velocidad, fidelidad o seguridad. La mayoría de los navegadores comparten una arquitectura de referencia compuesta por componentes esenciales para procesar y mostrar contenido web.
}

\nota{Para el examen}{
¿Qué factores determinan el proceso de interacción en un navegador?

\textbf{Respuesta:} El propósito y la configuración del navegador, como responder rápido, ser fiel al contenido o priorizar la seguridad de las comunicaciones.
}

\subtituloA{Proceso de interacción y dependencias}

\rista{
\lista{
  **Inicio y fin:** La interacción comienza con la indicación de la dirección del recurso y termina con su visualización. |
  **Dependencias:** El proceso se configura según el propósito (velocidad, fidelidad al contenido, seguridad de la comunicación con el servidor).
}
}
